\section{Unitary constructions}
\subsection{Established models and audited operations}
The baseline contains 84 realized input rings. Seventy-nine have selected
full tensor-model receipts; the Haagerup--Izumi $C_5$ model and four additional
Liptrap-positive inputs complete this count. Standard sources include
pointed and group-representation categories, $SU(2)$ and adjoint $PSU(2)$
categories, Tambara--Yamagami and metaplectic categories, Deligne products,
equivariantizations, group-theoretical categories, and explicit lower-rank
unitary models from Liu--Palcoux--Ren. Every selected receipt includes a
unit-preserving basis permutation and a complete coefficient comparison.

The construction audit checks more than tensor matching for the less standard
operations. It replays 27 associative $C_2$ zestings, eight central-extension
quotients, six cyclic actions on Tambara--Yamagami categories, and the
Ising-swap $C_4$ subextension. In a $C_2$ zesting, a neutral central
involution supplies the two-cocycle, and the fourth-cohomology obstruction
vanishes; the scalar coherence correction is included in the proof. Bases
for which the required central or braided data have not been supplied are
not used for this operation. The eight quotient constructions instead start
from $SU(2)_4$, a specified central extension of $V_4$, and a central boson;
their fusion tensors follow from the actual free-orbit quotient.

\subsection{Explicit coherent pointed actions}
Six new small models are equivariantizations of pointed unitary categories:
\begin{center}
\begin{tabular}{ll}
\toprule Inputs & Construction \\\midrule
$R9\_1\_47$ & $\Vec(C_{12},\omega)^{C_2}$, inversion action \\
$R9\_1\_98$ & $\Vec(C_3\times V_4,\omega)^{C_2}$, inversion and swap \\
$R10\_1\_61,84,109,158$ & $\Vec(\mathbb F_2^3,\omega)^{S_3}$, coordinate permutations \\
\bottomrule
\end{tabular}
\end{center}
For each model the certificate gives all associator exponents $w$, forward
tensorators $J_g$, and compositors $\mu_{g,h}$ modulo eight. The verifier
checks the pointed pentagon and, for every tuple of indices,
\begin{align*}
J_g(a,b)+J_g(ab,c)+w(a,b,c)
 &=w(ga,gb,gc)+J_g(b,c)+J_g(a,bc),\\
\mu_{g,h}(ab)+J_g(ha,hb)+J_h(a,b)
 &=J_{gh}(a,b)+\mu_{g,h}(a)+\mu_{g,h}(b),\\
\mu_{g,h}(ka)+\mu_{gh,k}(a)
 &=\mu_{h,k}(a)+\mu_{g,hk}(a).
\end{align*}
These equalities define coherent categorical actions. Independent code then
constructs all equivariant simple objects as projective matrix actions on
their homogeneous fibers, verifies the action equations, and computes every
fusion coefficient by exact Hom-space characters. Roots of unity of orders
sixteen or forty-eight suffice for these matrix calculations. Thus neither
the existence nor the fusion rules are inferred from a permutation action
on a based ring alone.

The dimension-160 cases $R10\_1\_39$ and $R10\_1\_138$ are likewise
realized as $D_{10}$-equivariantizations of $\Vec(\mathbb F_{16},\omega)$.
One may take $R10\_1\_39=\Rep(\mathbb F_{16}\rtimes D_{10})$, with
action $v\mapsto a^kv^{4^b}$. The second model uses a nontrivial associator
and explicitly stored fourth-root tensorators. The replay checks all 65,536
pointed pentagons, all action diagrams, the projective stabilizer data,
and the 1,000 fusion coefficients of each input.

\subsection{A twisted subgroup-bimodule model}
For $R9\_1\_92$, put $V=\{a\in\mathbb F_2^3:\sum_i a_i=0\}$ and
$G=V\rtimes S_3\cong S_4$. Restrict one of the coherent $S_3$ actions
above to $V$. The crossed-product associator is
\[
 \Omega((a,g),(b,h),(c,k))
 =w(a,gb,ghc)+J_g(b,hc)-\mu_{g,h}(c)\pmod8.
\]
It is verified directly on every quadruple of the 24-element group. Its
restriction to the chosen reflection subgroup $H=C_2$ is trivial. The
category of bimodules over $\bigoplus_{h\in H}h$ in $\Vec(G,\Omega)$
is therefore unitary. The supplied homogeneous bimodules are checked for
left, right, and mixed associativity; the independent reconstruction forms
their balanced tensor products, checks that both actions descend to the
quotient, and identifies all 729 coefficients with $R9\_1\_92$.

\subsection{The triality model of dimension twenty-four}
Nikshych--Riepel, Example 9.3, constructs a nontrivial $C_3$ extension
of $\Rep(D_8)$, integral of dimension 24, with every noninvertible simple
of dimension two. Their definition of nontrivial means that a homogeneous
component is not equivalent to the regular category as a left module.
An invertible simple in that component would give such an equivalence.
Thus both nonneutral components have exactly two dimension-two simples.

Write the neutral simples as $G=V_4$ and $m$, with $m^2=\sum G$.
Every rank-two invertible $\Rep(D_8)$-module has transitive $G$ action.
Indeed the module-category classification gives reflection subgroups $C_2$
or the nontrivial projective class of $D_8$; both are transitive under the
linear characters. The central $C_2$ alternative has a pointed dual and
cannot be an invertible $\Rep(D_8)$-bimodule. The other subgroups have
one, four, or five projective simples, so no case is omitted.

Label the nonneutral simples $X_0,X_1$ and $Y_0,Y_1$ with $X_i^*=Y_i$.
Their left and right actions are nonzero homomorphisms
$\phi,\psi:V_4\to C_2$. Multiplication by $m$ is the all-ones matrix.
Frobenius reciprocity determines the mixed products:
\[
 X_iY_j=\sum_{\phi(g)=i+j}g+m,\qquad
 Y_iX_j=\sum_{\psi(g)=i+j}g+m.
\]
The only unknowns are
$X_iX_j=a_{ij}Y_0+(2-a_{ij})Y_1$, with $a_{ij}\in\{0,1,2\}$;
duality determines $Y_jY_i$. The exact replay checks all
$3\cdot3\cdot3^4=729$ choices, allowing multiplicity two. Precisely six
are associative, all with distinct left and right kernels and $a_{ij}=1$.
Every one is based-isomorphic to $R9\_1\_117$. This identifies the
existence theorem's category with the supplied tensor. The category is
nilpotent and hence unitarizable by Galindo--Hong--Rowell, Theorem 5.20.

\subsection[An actual unitary model for R10-1-21]{An actual unitary model for $R10\_1\_21$}
Let $\mathcal P=PSU(2)_6$ have simples $1,f,x,y=fx$, with
$f^2=1$, $x^2=y^2=1+x+y$, and $xy=f+x+y$.
The invertible $f$ is a transparent fermion and its pointed associator is
trivial. All four simples are self-dual. Write
$\mathcal W=(\mathcal P\boxtimes\mathcal P)^{C_2}$ for the canonical
permutation equivariantization, a unitary braided fusion category.
Its simples are $d_i^\epsilon=(i\boxtimes i,\epsilon)$ and
$p_{ij}=(i\boxtimes j)\oplus(j\boxtimes i)$ for $i<j$.

The object $h=d_f^-$ is an order-two transparent boson: the two fermionic
self-braidings multiply to $+1$, its pointed associator is trivial, and
its equivariance squares to $+1$. Consequently $A=1\oplus h$ is a
connected central commutative separable algebra, and $\mathcal C=\mathcal W_A$
is an actual unitary fusion category. The free-module functor $F$ is tensor.
The $h$-orbits of size two give one simple each; a fixed simple splits
into two simples, since its stabilizer is cyclic of order two. There are
six free orbits and two fixed objects, namely $p_{1f}$ and $p_{xy}$.
Thus $\mathcal C$ has rank ten and dimensions
$1^4,(2+2\sqrt2)^2,(3+2\sqrt2)^4$.

\paragraph{The fixed invertibles, including their tensor squares.}
Put $s=F(d_1^-)$, and let $g,sg$ be the two summands of $F(p_{1f})$.
The pointed subcategory of $\mathcal P\boxtimes\mathcal P$ can be made
strict. For $X=p_{1f}$ write $e_1=1\boxtimes f$, $e_2=f\boxtimes1$.
The equivariant structure exchanges $e_1,e_2$, whereas that of $h$ has
sign $-1$. An equivariant right $h$-action on $X$ is therefore
\[
 r(e_1\otimes h)=i e_2,\qquad r(e_2\otimes h)=-i e_1;
\]
it squares to the identity. Its negative is the other simple module.
The half-braiding $c_{h,X}$ is $-1$ on both components, so the left
$h$-action used in the balanced product is $-r$. Hence $g\otimes_Ag$
is the cokernel of $r\otimes1+1\otimes r$ inside $X\otimes X$.
This operator is Hermitian, so that cokernel identifies with its kernel.
Its two homogeneous lines are
\[
 e_1e_1+e_2e_2,\qquad e_1e_2-e_2e_1.
\]
They have equivariance $+$ and $-$ respectively and underlie $1\oplus h$.
The remaining right action $1\otimes r$ sends the first line to $i$ times
the second and the second to $-i$ times the first; rescaling the second
line by $i$ identifies this with the regular right $A$-module.
Thus $g^2=1$. Multiplication by $s$ exchanges the two module structures,
so the four invertibles form $C_2\times C_2$, with $s^2=g^2=1$.
This calculation is essential: the choice $h=d_f^+$ gives a different
fixed-point resolution and does not give the multiplicity-one tensor here.

\paragraph{All fusion coefficients.}
Choose $a=F(p_{1x})$ and $b=F(d_x^+)$. Then
\[
\begin{array}{c|cccccccccc}
\text{source}&d_1^+&d_1^-&d_f^+&d_f^-&d_x^+&d_x^-&d_y^+&d_y^-&p_{1f}&p_{xy}\\
F&1&s&s&1&b&sb&sb&b&g+sg&gb+sgb
\end{array}
\]
and $F(p_{1y})=F(p_{fx})=ga$, $F(p_{fy})=a$.
Indeed $s$ interchanges the signs on the diagonal simples, and source
fusion gives $(g+sg)b=F(p_{xy})$ and $(g+sg)a=2F(p_{1y})$.
Since $sb\ne b$ and $g^2=1$, these identify a free orbit of four $b$'s
and an orbit of two $a$'s, with $sa=a$.
Ordinary equivariant tensor products give
\[
\begin{split}
b^2&=1+a+ga+b+sb+gb+sgb,\\
a^2&=1+s+a+ga+b+sb,\\
ab&=a+b+sb+gb+sgb.
\end{split}
\]
For example $d_x^+d_x^+=d_1^++d_x^++d_y^++p_{1x}+p_{1y}+p_{xy}$.
The three displayed rules, commutativity, and the invertible action
therefore determine the entire tensor without choosing any unknown
fusion coefficient. They also show all ten simples are self-dual.

The standalone constructor \nolinkurl{model21.py} generates these rules in the displayed model basis. It checks all $10^4$ associativity
identities, and verifies all $10^3$ coefficients against $R10\_1\_21$
under the zero-based model-to-input basis map
\[
(0,9,1,2,3,8,6,7,5,4).
\]
The resulting model receipt is \nolinkurl{model21.json}. Its inverse
input-to-model map is $(0,2,3,4,9,8,6,7,5,1)$. The independent standard-library
verifier \nolinkurl{work/audit_new_models/verify_model21.py} replays the
fixed-point balance, all $196$ source products, the free-module map, and
the complete original-tensor comparison. Its receipt is
\nolinkurl{work/audit_new_models/model21_audit.json}.
The categorical inputs are the standard unitary $PSU(2)_6$ category,
canonical permutation equivariantization, and central Tannakian
condensation (EGNO, Theorems 8.4.9 and 8.23.3). No reconstruction from
a ring-consistent induction matrix is being asserted.

\subsection[Triality equivariantizations]{Triality equivariantizations: $R9\_1\_12$ and $R9\_1\_108$}
\label{sec:triality-equivariant}

\paragraph{Audit status.}
Both constructions below have passed independent categorical and
coefficient audits. The second uses a direct inner modification of the
first action and the vanishing of its action-coherence obstruction.

Put $D=\operatorname{Rep}(D_8)$, where $D_8$ has order eight.  We use its
Tambara--Yamagami presentation
$D=\operatorname{TY}(V_4,\chi,+\tfrac12)$, with alternating bicharacter
$\chi$.  The standard associators are invariant under every permutation of
the three nonidentity elements of $V_4$.  Consequently these permutations,
with identity tensorators, define a coherent unitary $S_3$ action.
The tensor autoequivalence group is $S_3$, while
$\operatorname{BrPic}(D)\cong S_4$; see Nikshych--Riepel,
\emph{Categorical Lagrangian Grassmannians and Brauer--Picard groups of
pointed fusion categories}, \href{https://arxiv.org/abs/1309.5026}{arXiv:1309.5026},
Section~9.1, and Diatlyk--Luo--Wang--Weller,
\emph{Gauging Non-Invertible Symmetries: Topological Interfaces and
Generalized Orbifold Groupoid in 2d QFT},
\href{https://arxiv.org/abs/2311.17044}{arXiv:2311.17044}, Section~3.5,
Equation~(3.30), for the explicit tensor triality.

\paragraph{A coherent Morita transport.}
The three nonregular invertible left $D$-module classes of rank two are the
two reflection-subgroup modules and the projective $D_8$-representation
module for the nontrivial Schur class.  Their simple objects have internal
End objects $1+h+m$, with the three distinct nontrivial choices of $h$.
Tensor triality therefore permutes these three module classes transitively.
Choose the projective module class and realize it by the Azumaya algebra
$A=\operatorname{End}(V)$ in the symmetric category $D$, where $V$ is a real
two-dimensional irreducible projective representation of $D_8$ arising
from its central cover $D_{16}$.  Thus $A=1+h+m$ and
$A\text{-}\operatorname{Bimod}(D)\simeq D$.

Choose the transposition represented by the actual group involution
\[
 a(r)=r^{-1},\qquad a(s)=sr,
 \qquad r^4=s^2=1,\quad srs=r^{-1}.
\]
It fixes the projective module class.  This invariance has a coherent lift:
the nontrivial Schur cocycle of $D_8$ extends to
$D_8\rtimes_a C_2\cong D_{16}$, because the cover $D_{32}\to D_{16}$
restricts to the nontrivial cover $D_{16}\to D_8$ on its normal $D_8$
subgroup.  Equivalently, writing the projective representation as a
rotation $R$ through $\pi/4$ and a reflection $S$, the real reflection
$T=SR^{1/2}$ satisfies
\[
 T^2=1,\qquad TRT=R^{-1},\qquad TST=SR,
\]
and conjugation by $T$ makes $A$ a coherent equivariant algebra.

There is also no ambiguity between this coherent group action and the
strict TY transposition.  A fixed transposition class has two coherent
$C_2$ actions, distinguished by
$H^2(C_2,\operatorname{Aut}_\otimes(\operatorname{Id}_D))
=H^2(C_2,C_2)=C_2$.
For the strict action the two extensions of $m$ have equivariant scalars
$+1,-1$ and are self-dual; in the other class they have scalars $+i,-i$
and form a dual pair.  The group-automorphism action has self-dual
extensions in $\operatorname{Rep}(D_{16})$, so it is coherently equivalent
to the strict restriction.

Form the crossed product $B=D\rtimes S_3$ and its graded Morita dual
$C=A\text{-}\operatorname{Bimod}(B)$.  The neutral component is again $D$.
Its $C_3$ restriction $E$ is nonregular, since triality moves the module
class of $A$; the preceding triality construction identifies its fusion
ring with $R9\_1\_117$.  Coherent $C_2$ equivariance of $A$ supplies an
invertible $t\in C$ with $t^2\cong1$ and trivial restricted associator.
Thus conjugation by $t$ is an actual coherent $C_2$ action on $E$, and
$C=E\rtimes C_2$.  On the neutral category its equivariantization is
$\operatorname{Rep}(D_{16})$.  Indeed, if $q=st$ in
$D_8\rtimes_a C_2$, then $q^2=r$ and $sqs=q^{-1}$.

\paragraph{Duality of the dimension-four objects.}
The subcategories of a faithfully graded fusion category supported on
conjugate subgroups are Morita equivalent: the sum of components on the
corresponding coset is the equivalence bimodule.  Hence all three
order-two subgroup subextensions of $C$ are Morita equivalent to
$D\rtimes C_2$, and therefore to $\operatorname{Rep}(D_{16})$.
The two nonregular ones have dimension profile $(1^4,2^3)$.
Every category in this Morita class is
$\mathcal C(D_{16},1,H,\psi)$.  This profile forces $H=D_{16}$ or a
reflection $C_2$.  Cyclic normal subgroups give pointed duals; a
Klein-four subgroup or an index-two $D_8$ subgroup already gives eight
invertibles, for either available Schur class.  For $H=D_{16}$ the fusion
ring is the representation ring of $D_{16}$.  For a reflection $C_2$,
the three non-normalizer double cosets are inverse-stable with trivial
stabilizer, and the four invertibles form $V_4$.  In both permitted
cases all simples are self-dual.

Write the nonneutral $C_3$ components of $E$ as
$\{X_0,X_1\}$ and $\{Y_0,Y_1\}$, with $X_i^*=Y_i$.
The self-duality just proved for $X_it$ implies $t(X_i)=Y_i$.
Thus both four-dimensional orbit objects in $E^{C_2}$ are self-dual.

\paragraph{Exact tensor identification.}
Label $E$ by
$0,1,2,3,m,X_0,X_1,Y_0,Y_1$, choosing left and right stabilizers
$\{0,1\}$ and $\{0,2\}$ for the $X$ component.
The restriction of the seven neutral simples of
$\operatorname{Rep}(D_{16})$ is
\[
 0,\ 3,\ 0,\ 3,\ m,\ 1+2,\ m,
\]
and the other two restrictions are $X_i+Y_i$.
Let $F$ be this restriction matrix.  Induction is $F^{\mathsf T}$ by
Frobenius reciprocity, and
\[
 \operatorname{Ind}(X)\otimes W
 =\operatorname{Ind}(X\otimes\operatorname{Forget}(W)),\qquad
 W\otimes\operatorname{Ind}(X)
 =\operatorname{Ind}(\operatorname{Forget}(W)\otimes X).
\]
These formulas determine every product outside the neutral subcategory;
inside it we use the ordinary $D_{16}$ representation ring.  The resulting
multiplicity-free tensor agrees with $R9\_1\_12$ under
\[
 (0,3,2,4,8,7,6,1,5).
\]
The independent verifier
\nolinkurl{work/audit_new_models/verify_triality_equivariant.py}
checks all $729$ coefficients, the restriction homomorphism, and the
$D_{16}$ subgroup calculation.  The certificate is
\nolinkurl{work/resolve_small_models/triality_equivariant_certificate.json}.
The construction is unitary throughout; alternatively, unitarizability
follows from Galindo--Hong--Rowell,
\href{https://arxiv.org/abs/1105.5048}{arXiv:1105.5048}, Theorem~5.20.

\paragraph[The second coherent action and R9-1-108]{The second coherent action and $R9\_1\_108$.}
The following direct construction has passed the independent categorical
audit and supplies the second unitary realization.
Let $\alpha$ denote the accepted action and choose $g\in V_4$ with
$g\alpha(g)=h$, the nontrivial fixed invertible.  Set
$\beta=\operatorname{Ad}_g\circ\alpha$.  Then
$\beta^2\simeq\operatorname{Ad}_h$.
For the nonregular triality action, the fixed central invertible in
$\operatorname{Inv}Z(D)$ is $w=(z,h)$, where $z$ is the central involution
of $D_8$.  Its obstruction to a central lift on $h\in E$ lies in
$H^2(C_3,\mathbb C^\times)=0$; hence $\operatorname{Ad}_h\simeq\operatorname{Id}$.
The obstruction to making $\beta$ a coherent $C_2$ action lies in
\[
 H^3(C_2,\operatorname{Aut}_\otimes(\operatorname{Id}_E))
 =H^3(C_2,C_3)=0,
\]
and the corresponding $H^2$ torsor is trivial.  Any correction is a
character of $U(E)=C_3$, so it is the identity on $D$.
Relative to the symmetric central lift of $h$, the lift $w$ differs by
$z$, changing the neutral square coherence by $-1$ on $m$.  The inner
modification changes the orbit pairing to $X_i+Y_{i+1}$.
These obstruction statements use Etingof--Nikshych--Ostrik,
\emph{Fusion categories and homotopy theory},
\href{https://arxiv.org/abs/0909.3140}{arXiv:0909.3140},
Corollary~7.9(i)~\cite{ENOHomotopy},
and Gelaki--Naidu--Nikshych,
\emph{Centers of graded fusion categories},
\href{https://arxiv.org/abs/0905.3117}{arXiv:0905.3117}, Theorem~3.5.
The detailed central-lift and coherence comparison is in
\nolinkurl{work/resolve_integral/TRIALITY_ACTIONS.md}.

The neutral ring is therefore $\operatorname{Rep}(SD_{16})$, since the
corresponding extension has $t^2=z=r^2$, and $q=st$ satisfies
$q^8=s^2=1$ and $sqs=q^3$.
Replacing the neutral group ring accordingly and using this
pairing in the same induction formulas gives exactly $R9\_1\_108$,
with based permutation
\[
 (0,4,6,8,7,1,2,3,5).
\]
Both this coefficient identification and the categorical coherence of the
modification have passed independent audits. No classification of all
possible involutions is required for these two positive witnesses.


